\scanpage{49}
% archive sheet number: 48
\mspage{15}
topologiques paracompacts, à un théorème d'isomorphisme pour les \deleted{préfa} préfaisceaux \handwritten{[}T.F\quad Chap. II, \S~5, Th. 10.2\handwritten{]} comme on s'en convaint facilement en examinant le cas des préfaisceaux constants.

% slanted
\margin{\handwritten{aligner}}
Cependant, la notion introduite ici coincide avec celle de \handwritten{[}T.F.\handwritten{]} dans le cas particulier des préfaisceau abéliens $P$ tel que $P(\emptyset)=0$, condition vérifiée par exemple par les préfaisceaux $\handwritten{\mathscr H}^q(F)$ \deleted{\unsure{la|le|ta}} qui s'introduisent dans la suite spectrale de Leray.

Enfin soit \deleted{\unsure{x|a|A}} $a:C_A^{\handwritten{\wedge}}\handwritten{\longrightarrow}C_A^{\handwritten{\sim}}$ le foncteur faisceau associé. Le foncteur $a$ est exact. Il est adjoint à gauche au foncteur $\handwritten{\mathscr H^0}$. Le morphisme d'adjonction :
\[
a\handwritten{\mathscr H^0\longrightarrow}\operatorname{id}_{C_A^{\handwritten{\sim}}}\tag*{(21)}
\]
est un isomorphisme (2.3.14.II).

La formule d'adjonction fournit pour tout préfaisceau d'ensembles $S$ sur $C$ et tout faisceau d\deleted{e}\handwritten{'} $A$-modules $F$ un isomorphisme :

\inserted{\handwritten{$F(aS)\xrightarrow{\sim}\mathscr H^0(F)(S)$ i.e. :}}
\[
\deleted{H^0(\unsure{C|C^{\sim}}/aS,F)\ \unsure{\longrightarrow|=}\ H^0(C^{\unsure{\wedge|\sim}}/S,F)\unsure{=|-}H^0(C/S,}
\]
\[
H^0(C^{\handwritten{\sim}}/aS,F)\handwritten{\xrightarrow{\sim}}H^0(C^{\handwritten{\wedge}}/S,\handwritten{\mathscr H^0}(F))\tag*{(22)}
\]

De plus si $G$ est un préfaisceau \unsure{dé|de} $A$-modules tel que le faisceau associé soit nul, on a ((2.3.12.II) et (2.3.14.II)) :
\[
\handwritten{\check{\mathscr H}^{0}}(G)=0\tag*{(23)}
\]

\margin{\handwritten{Laisser ici un interligne $\longmapsto$}}
Nous allons maintenant én\deleted{n}oncer les propositions liant la cohomologie de $\overset{\handwritten{\scriptscriptstyle\vee}}{\text{C}}$ech à la cohomologie des faisceaux.

% Hand-drawn left bracket spans the proposition.
\uline{Proposition 2.2} : Le foncteur $\handwritten{\mathscr H^0}$ transforme les faisceaux injectifs en préfaisceaux injectifs. Soit $F$ un faisceau de $A$-modules. Les faisceaux associés aux préfaisceaux $\handwritten{\mathscr H^q}(F)$ $(q>0)$ sont nuls.

\uline{Preuve} : Le foncteur faisceau associé , adjoint à gauche du foncteur $\handwritten{\mathscr H^{\unsure{0|c}}}$ est exact ; d'où la première assertion (0.2). Le foncteur $a\,\handwritten{\mathscr H^0}$ est isomorphe au foncteur identique (Formul\deleted{\unsure{l|i}}\unsure{l|{}}e (21)) d'où la deuxième assertion.

% Hand-drawn left bracket spans the proposition and its display.
\uline{Proposition 2.3} : Soit $S$ un préfaisceau d'ensembles sur $C$.

Il existe une suite spectrale fonctoriel\handwritten{\inserted{le}} en $S$ et $F$ :
\[
\deleted{H^{p+q}(\unsure{C|C^\sim}/\unsure{U|u},F)\unsure{\Leftarrow|\mathbin{\text{---}}|{}}\unsure{H|H^p}}
\]
\[
H^{p+q}(C^{\handwritten{\sim}}/aS,F)\handwritten{\Longleftarrow}H^p(C^{\handwritten{\wedge}}/S,\handwritten{\mathscr H^q}(\unsure{F|E}))
\]
